% Źródła: Eves s. 287, 306--307, 352 (zad. 6--7); Wikipedia: Eugenio Beltrami, Arthur Cayley, Felix Klein, Beltrami–Klein model, Non-Euclidean geometry.
% https://en.wikipedia.org/wiki/Beltrami%E2%80%93Klein_model
% na podstawie https://en.wikipedia.org/wiki/Beltrami–Klein_model
% model Kleina - Delta 2012-06

\section{Beltrami, Cayley i Klein} % lang-pl
\section{Beltrami, Cayley e Klein} % lang-it
\label{sekcja_beltrami_cayley_klein}

Trzej matematycy, którzy prawie się nie spotkają, zbudują model, w którym płaszczyzna Łobaczewskiego stanie się wnętrzem koła. % lang-pl
W tej sekcji zobaczymy, kto co zrobi i kto komu odpowie -- oraz dlaczego model dostanie nazwę od nazwiska tego, który nie będzie pierwszy. % lang-pl
Tre matematici che quasi non s'incontreranno costruiranno un modello in cui il piano di Lobachevskij diventerà l'interno di un cerchio. % lang-it
In questa sezione vedremo chi farà che cosa e chi risponderà a chi -- e perché il modello prenderà il nome di colui che non sarà il primo. % lang-it

Eugenio Beltrami (1835--1900) urodzi się w Cremonie, będzie studiował w Pawii, potem pracował w kolejach, a od 1862 roku wykładał w Bolonii, następnie w Pizie, Rzymie i znów w Pawii, a od 1891 roku w Rzymie. % lang-pl
Eugenio Beltrami (1835--1900) nascerà a Cremona, studierà a Pavia, lavorerà poi per le ferrovie e dal 1862 insegnerà a Bologna, quindi a Pisa, a Roma e di nuovo a Pavia, e dal 1891 a Roma. % lang-it
\index[persons]{Beltrami, Eugenio}%

W 1868 roku opublikuje dwie prace po włosku (w 1869 roku Jules Hoüel wyda ich tłumaczenie francuskie): \emph{Saggio di interpretazione della geometria non-euclidea} i \emph{Teoria fondamentale degli spazii di curvatura costante}. % lang-pl
W pierwszej pokaże, że geometria nieeuklidesowa realizuje się na powierzchni zwanej pseudosferą (jej geodezyjne badali wcześniej Ferdinand Minding w 1840 roku i Delfino Codazzi w 1857 roku), a w drugiej przedstawi ogólne rozumowanie o niesprzeczności w dowolnej liczbie wymiarów; powoła się przy tym na wykład Riemanna z 1854 roku (zobacz sekcję \ref{sekcja_riemann}). % lang-pl
Nel 1868 pubblicherà due memorie in italiano (nel 1869 Jules Hoüel ne darà la traduzione francese): \emph{Saggio di interpretazione della geometria non-euclidea} e \emph{Teoria fondamentale degli spazii di curvatura costante}. % lang-it
Nella prima mostrerà che la geometria non euclidea si realizza su una superficie detta pseudosfera (le cui geodetiche avevano studiato prima Ferdinand Minding nel 1840 e Delfino Codazzi nel 1857), nella seconda presenterà un ragionamento generale di coerenza in un numero qualsiasi di dimensioni; si richiamerà a tal fine alla conferenza di Riemann del 1854 (si veda la sezione \ref{sekcja_riemann}). % lang-it
\index[persons]{Hoüel, Jules}%
\index[persons]{Minding, Ferdinand}%
\index[persons]{Codazzi, Delfino}%

Kolejnym osiągnięciem będzie model Beltramiego-Kleina dla geometrii hiperbolicznej, w którym pracujemy z punktami wewnątrz otwartego koła. % lang-pl
Proste w tym modelu to cięciwy łączące punkty z brzegu koła. % lang-pl
Wspomni o nim Eugenio Beltrami w 1868 roku, najpierw dla $n = 2$, a potem dla dowolnej liczby wymiarów. % lang-pl
Un ulteriore progresso sarà il modello di Beltrami-Klein della geometria iperbolica, nel quale si lavora con i punti all'interno di un disco aperto. % lang-it
Le rette di questo modello sono le corde che uniscono punti del bordo del disco. % lang-it
Eugenio Beltrami lo presenterà nel 1868, dapprima per $n=2$ e successivamente per un numero arbitrario di dimensioni. % lang-it
Jego prace dowiodą, że geometria hiperboliczna jest równoniesprzeczna z euklidesową. % lang-pl
Le sue opere proveranno che la geometria iperbolica è equiconsistente con quella euclidea. % lang-it

% memoirs
Wcześniej, bo w 1859 roku, Arthur Cayley użyje dwustosunkowej definicji kąta (od Laguerre'a) by pokazać, że geometria euklidesowa może być zdefiniowana przy pomocy geometrii rzutowej. % lang-pl
Arthur Cayley urodzi się 16 sierpnia 1821 roku w Richmond pod Londynem, pierwsze osiem lat życia spędzi w Petersburgu, a umrze 26 stycznia 1895 roku w Cambridge; jego przyjacielem będzie James Sylvester, z którym będą spacerować po dziedzińcach Lincoln's Inn, rozmawiając o niezmiennikach. % lang-pl
In precedenza, nel 1859, Arthur Cayley utilizzerà la definizione di angolo tramite rapporto anarmonico (dovuta a Laguerre) per mostrare che la geometria euclidea può essere definita mediante la geometria proiettiva. % lang-it
Arthur Cayley nascerà il 16 agosto 1821 a Richmond, presso Londra, trascorrerà a San Pietroburgo i primi otto anni di vita e morirà il 26 gennaio 1895 a Cambridge; suo amico sarà James Sylvester, con il quale passeggerà nei cortili di Lincoln's Inn discutendo di invarianti. % lang-it
\index[persons]{Cayley, Arthur}%
\index[persons]{Laguerre, Edmond}%
\index[persons]{Sylvester, James}%

Ten wynik ukaże się w \emph{A sixth memoir upon quantics}; odległość, którą tam zdefiniuje za pomocą logarytmu dwustosunku, zostanie później nazwana metryką Cayleya. % lang-pl
Cayley nazwie ponadto zbiór punktów idealnych płaszczyzny Łobaczewskiego jej \emph{absolutem} \cite[s. 307]{eves1_1972} (zobacz sekcję \ref{sekcja_punkty_idealne}). % lang-pl
Questo risultato apparirà in \emph{A sixth memoir upon quantics}; la distanza che vi definirà tramite il logaritmo del birapporto sarà poi chiamata metrica di Cayley. % lang-it
Cayley chiamerà inoltre il luogo dei punti ideali del piano di Lobachevskij il suo \emph{assoluto} \cite[p. 307]{eves1_1972} (si veda la sezione \ref{sekcja_punkty_idealne}). % lang-it

Felix Klein (1849--1925) urodzi się w Düsseldorfie i umrze w Getyndze. % lang-pl
Klein zapozna się z pracami Beltramiego i Cayleya; w 1870 roku, jako dwudziestolatek, weźmie udział w seminarium Weierstrassa, gdzie omówi pomysły Cayleya i zapyta: czy między pomysłami Cayleya oraz Łobaczewskiego istnieje powiązanie? % lang-pl
Usłyszy w odpowiedzi, że oba systemy dzieli przepaść pojęciowa. % lang-pl
Felix Klein (1849--1925) nascerà a Düsseldorf e morirà a Gottinga. % lang-it
Felix Klein studierà i lavori di Beltrami e Cayley; nel 1870, appena ventenne, parteciperà al seminario di Weierstrass, dove discuterà le idee di Cayley e chiederà: esiste una relazione tra le idee di Cayley e quelle di Lobachevskij? % lang-it
Riceverà in risposta che i due sistemi sono concettualmente molto lontani. % lang-it
\index[persons]{Klein, Felix}%
\index[persons]{Weierstrass, Karl}%

Klein nie zrezygnuje: w 1871 roku opublikuje pracę \emph{Über die sogenannte Nicht-Euklidische Geometrie} i pokaże w niej, że metody Cayleya dają rzutowy model geometrii nieeuklidesowej. % lang-pl
Wprowadzi wtedy nazwy \emph{hiperboliczna} i \emph{eliptyczna} (geometrię euklidesową nazwie paraboliczną, ale ta nazwa nie utrwali się), a rok później, w 1872, wygłosi program erlangeński. % lang-pl
Choć Beltrami był pierwszy, model przejdzie do historii jako model Kleina (a sprawiedliwiej -- Cayleya--Kleina); prace Beltramiego z 1868 roku pozostaną w dużej mierze niezauważone. % lang-pl
Klein non si arrenderà: nel 1871 pubblicherà il lavoro \emph{Über die sogenannte Nicht-Euklidische Geometrie} e vi mostrerà che i metodi di Cayley forniscono un modello proiettivo della geometria non euclidea. % lang-it
Introdurrà allora i nomi \emph{iperbolica} ed \emph{ellittica} (chiamerà parabolica la geometria euclidea, ma questo nome non si affermerà), e un anno dopo, nel 1872, terrà il programma di Erlangen. % lang-it
Anche se Beltrami era stato il primo, il modello passerà alla storia come modello di Klein (più giustamente: di Cayley--Klein); i lavori di Beltrami del 1868 rimarranno in gran parte inosservati. % lang-it
\index{model!Beltramiego-Kleina} % lang-pl
\index{modello!di Beltrami-Klein} % lang-it

\begin{definition}
[model Beltramiego-Kleina] % lang-pl
[modello di Beltrami-Klein] % lang-it
	Ustalmy okrąg $\Sigma$. % lang-pl
	Punktami są punkty wnętrza $\Sigma$, prostymi -- cięciwy $\Sigma$, a długość odcinka $AB$ to $\log(AB, TS)$, gdzie $S$, $T$ są końcami cięciwy zawierającej $AB$, uporządkowanymi tak, że $A$ leży między $S$ i $B$. % lang-pl
	Odcinki (kąty) są przystające, gdy jeden przechodzi na drugi przez rzutowość przeprowadzającą $\Sigma$ i jej wnętrze na siebie. % lang-pl
	Fissiamo una circonferenza $\Sigma$. % lang-it
	I punti sono i punti dell'interno di $\Sigma$, le rette sono le corde di $\Sigma$, e la lunghezza del segmento $AB$ è $\log(AB, TS)$, dove $S$, $T$ sono gli estremi della corda contenente $AB$, ordinati in modo che $A$ stia tra $S$ e $B$. % lang-it
	Segmenti (angoli) sono congruenti quando uno si trasforma nell'altro tramite una proiettività che manda $\Sigma$ e il suo interno in sé. % lang-it
\end{definition}

Eves \cite[s. 352, zad. 6--7]{eves1_1972}, gdzie przypisuje ten model Kleinowi. % lang-pl
Eves \cite[p. 352, es. 6--7]{eves1_1972}, dove attribuisce questo modello a Klein. % lang-it

Beltrami odegra jeszcze jedną rolę w naszej opowieści: w 1889 roku przypomni światu zapomnianą książkę Saccheriego (zobacz sekcję \ref{sekcja_saccheri}). % lang-pl
Beltrami avrà ancora un altro ruolo nella nostra storia: nel 1889 richiamerà l'attenzione del mondo sul dimenticato libro di Saccheri (si veda la sezione \ref{sekcja_saccheri}). % lang-it
